\documentclass[11pt,a4paper]{article}

\usepackage[utf8]{inputenc}
\usepackage[margin=1in]{geometry}
\usepackage{amsmath,amssymb,amsfonts,mathrsfs}
\usepackage{graphicx}
\usepackage{booktabs}
\usepackage{multirow}
\usepackage{array}
\usepackage{xcolor}
\usepackage{hyperref}
\usepackage{cite}
\usepackage{float}
\usepackage{caption}
\usepackage{subcaption}
\usepackage{url}
\usepackage{microtype}
\usepackage{enumitem}

\hypersetup{
    colorlinks=true,
    linkcolor=blue!70!black,
    citecolor=blue!70!black,
    urlcolor=blue!70!black
}

\title{\textbf{Semantically Intelligent Epilepsy Surveillance: An Ontology-Driven Electronic Health Record Framework Integrating MIMIC-IV v3.1, HL7 FHIR R4, and Machine Learning for Population-Scale Neurological Analytics}}

\author{
  \textbf{M J Yogesh}$^1$ \quad \textbf{Dr. Karthikeyan J}$^{1,*}$ \\[0.5em]
  $^1$School of Computer Science Engineering and Information Systems,\\
  Vellore Institute of Technology, Vellore 632 014, Tamil Nadu, India \\[0.3em]
  \small{$^*$Corresponding Authors: \href{mailto:yogeshmj.nie@gmail.com}{yogeshmj.nie@gmail.com}, \href{mailto:karthikeyan.jk@vit.ac.in}{karthikeyan.jk@vit.ac.in}}
}

\date{\today}

\begin{document}

\maketitle

\begin{abstract}
Epilepsy affects over 50 million individuals worldwide, yet existing electronic health record (EHR) systems lack the semantic precision, ontological structure, and analytical depth required for population-level surveillance and clinical decision support. We present \textbf{EHR-EpSO}, a public health informatics framework addressing three persistent structural gaps---semantic heterogeneity, poor interoperability, and limited predictive capacity---by unifying the Epilepsy and Seizure Ontology (EpSO) with HL7 Fast Healthcare Interoperability Resources Release 4 (FHIR R4) standards and a modular machine learning analytics engine. Validated on the largest publicly available epilepsy cohort drawn from MIMIC-IV v3.1 (4,187 admissions, 3,641 unique patients spanning 2008--2022), our three-stage ontology alignment pipeline (prefix-tree lookup over 3,847 EpSO concept labels coupled with BioWordVec contextual disambiguation) achieved a concept-resolution rate (CR) of 94.1\% (95\% CI: 93.2\%--95.0\%), significantly outperforming general-purpose clinical NLP tools (MetaMap: 81.4\%, QuickUMLS: 84.7\%, scispaCy: 86.2\%). Furthermore, 99.2\% of admissions generated fully validated FHIR R4 resource bundles (\texttt{Condition}, \texttt{Observation}, \texttt{MedicationRequest}). Evaluated under a rigorous 5-fold stratified cross-validation protocol combined with an independent held-out test evaluation ($n=629$), a Transformer-encoder classifier achieved an AUROC of 0.876 (95\% CI: 0.856--0.896) for 30-day seizure recurrence and 0.839 (95\% CI: 0.815--0.863) for treatment non-adherence, significantly outperforming XGBoost and BiLSTM baselines ($p < 0.01$, DeLong's test). Dual-model interpretability using TreeSHAP for XGBoost and Integrated Gradients for the Transformer revealed prior admission frequency and Medicaid coverage as primary predictors of recurrence and non-adherence, respectively. Systematic ablation studies confirmed that every pipeline stage contributes significantly to performance. EHR-EpSO offers an open-source, production-scalable foundation for neurological population health analytics.

\vspace{0.5em}
\noindent \textbf{Keywords:} Epilepsy; Public Health Informatics; Electronic Health Records; Ontology; EpSO; HL7 FHIR R4; MIMIC-IV v3.1; Seizure Recurrence; Machine Learning; Interpretability.
\end{abstract}

\clearpage

\section{Introduction}

Epilepsy is one of the most common serious neurological conditions globally, affecting over 50 million individuals and contributing substantially to global disease burden \cite{who2019epilepsy, fisher2014ilae}. Low- and middle-income countries bear a disproportionate burden, with treatment gaps exceeding 75\% due to diagnostic delays, resource constraints, and fragmented health infrastructure \cite{who2019epilepsy}. Even in high-income healthcare settings, neurological care faces severe challenges: approximately 30\% of epilepsy patients develop drug-resistant epilepsy unresponsive to two or more appropriately chosen antiseizure medications (ASMs) \cite{kwan2000early, kwan2011definition}, while medication non-adherence rates range between 29\% and 58\%, driving avoidable status epilepticus admissions, emergency visits, and elevated mortality \cite{hovinga2008association}. Addressing these challenges requires population-scale clinical surveillance capable of converting unstructured Electronic Health Record (EHR) data into actionable clinical insights.

The proliferation of EHR systems and large-scale open clinical databases, such as MIMIC-IV v3.1 \cite{johnson2024mimic}, provides unprecedented opportunities for data-driven neurological research. Spanning over 364,000 hospital admissions from 2008 to 2022 at Beth Israel Deaconess Medical Center, MIMIC-IV offers rich longitudinal data comprising diagnostic codes, laboratory panels, pharmacy administration logs, and free-text clinical notes. However, translating this data wealth into effective clinical decision support and population surveillance is hindered by three persistent structural gaps:

\begin{enumerate}[label=\textbf{Gap \arabic*:}, leftmargin=*, itemsep=0.3em]
    \item \textbf{Semantic Heterogeneity in Clinical Narratives:} Epilepsy documentation exhibits extreme terminology variation. Clinical notes feature unstandardized abbreviations, non-uniform seizure descriptors, and inconsistent diagnostic codings that lag behind modern International League Against Epilepsy (ILAE) operational classifications \cite{fisher2017operational}. Standard vocabulary systems like SNOMED CT and ICD-10 lack the domain-specific granularity required to represent complex epilepsy phenotypes, seizure etiologies, and electroencephalographic (EEG) patterns.
    \item \textbf{Interoperability and Data Exchange Deficits:} Most existing clinical NLP and epilepsy predictive models operate on proprietary, siloed database schemas. The absence of standards-compliant serializations, such as HL7 Fast Healthcare Interoperability Resources (FHIR R4) \cite{hl7fhirr4}, prevents seamless data exchange, multi-center federated evaluation, and real-time clinical decision support integration.
    \item \textbf{Limited Predictive Precision and Black-Box Models:} Existing risk-stratification models for epilepsy rely predominantly on static logistic regressions or unvalidated rule-based heuristics. Modern deep learning architectures (e.g., Transformers, Recurrent Neural Networks) and ensemble methods hold strong predictive potential \cite{rajkomar2018scalable, vaswani2017attention}, but their adoption in neurology remains limited by small sample sizes, flawed evaluation methodologies, lack of formal model interpretability, and missing ablation validation.
\end{enumerate}

To bridge these gaps, we present \textbf{EHR-EpSO}, an open-source, ontology-driven public health informatics framework. EHR-EpSO integrates the Epilepsy and Seizure Ontology (EpSO) \cite{buccella2011ontology, larson2013neurolex} with HL7 FHIR R4 data standards and a modular machine learning engine. Validated on a cohort of 4,187 hospital admissions from 3,641 patients in MIMIC-IV v3.1, EHR-EpSO systematically addresses the limitations identified in prior literature.

\subsection{Summary of Main Contributions}
The key contributions of this work are as follows:

\begin{itemize}[leftmargin=*, itemsep=0.2em]
    \item \textbf{First Scalable FHIR R4-Compliant Epilepsy Framework:} We introduce an automated pipeline that ingests unstructured EHR narratives, maps them to the EpSO domain ontology, and serializes clinical entries into validated HL7 FHIR R4 resource bundles (\texttt{Condition}, \texttt{Observation}, \texttt{MedicationRequest}) with a 99.2\% compliance rate.
    \item \textbf{Hybrid NLP Alignment Pipeline with Comparative Benchmarking:} We design a three-stage ontology alignment architecture combining customized text normalization, prefix-tree (trie) indexing over 3,847 EpSO concept labels, and 300-dimensional BioWordVec contextual embedding disambiguation \cite{zhang2019biowordvec}. We demonstrate a concept-resolution rate of 94.1\%, outperforming standard clinical NLP baselines including MetaMap \cite{aronson2001overview}, QuickUMLS \cite{soldaini2016quickumls}, and scispaCy \cite{neumann2019scispacy}.
    \item \textbf{Methodologically Harmonized Machine Learning Evaluation:} Resolving evaluation protocol ambiguities common in prior work, we evaluate predictive models under a combined 5-fold stratified cross-validation and independent held-out test split ($n=629$). A Transformer encoder architecture achieves superior discrimination (AUROC 0.876 for 30-day seizure recurrence; 0.839 for medication non-adherence) compared to XGBoost \cite{chen2016xgboost} and BiLSTM \cite{hochreiter1997long} baselines ($p < 0.01$, DeLong's test).
    \item \textbf{Dual-Model Interpretability and Clinical Attribution:} We present comprehensive model interpretability using TreeSHAP for XGBoost \cite{lundberg2017unified} and Integrated Gradients for the Transformer \cite{sundararajan2017axiomatic}. We identify prior 12-month admission frequency and Medicaid insurance coverage as key actionable drivers of seizure recurrence and non-adherence, respectively.
    \item \textbf{Extensive Sensitivity Analysis and Systematic Ablations:} We provide empirical sensitivity sweeps over the vector disambiguation threshold $\tau \in [0.50, 0.80]$ and conduct rigorous ablation studies confirming the incremental contribution of each pipeline component and feature set.
    \item \textbf{Open-Source, Scalable Implementation:} Framework code, preprocessing scripts, FHIR conversion modules, and database adapters (supporting SQLite for local prototyping and PostgreSQL for hospital-scale concurrent throughput) are made publicly available at \url{https://github.com/yogeshmj2024/EHR-EpSO}.
\end{itemize}

\section{Related Work}

\subsection{Clinical Ontologies and Biomedical NLP}
Ontology-driven data integration has become indispensable for biomedical knowledge representation. The Unified Medical Language System (UMLS) Metathesaurus \cite{bodenreider2004unified} and SNOMED CT provide broad domain coverage but lack the specialized multi-dimensional taxonomy necessary for epilepsy, such as the ILAE 2017 seizure classification framework \cite{fisher2017operational}. The Epilepsy and Seizure Ontology (EpSO) \cite{buccella2011ontology}, developed within the NeuroLex framework \cite{larson2013neurolex}, models epilepsy concepts across four explicit dimensions: seizure semiology, anatomical localization, etiology, and comorbid neurological conditions. 

Traditional clinical NLP pipelines rely on lexical lookup engines such as MetaMap \cite{aronson2001overview} or string-matching systems like QuickUMLS \cite{soldaini2016quickumls}. More recent toolkits like scispaCy \cite{neumann2019scispacy} leverage specialized spaCy models for biomedical entity extraction. In contextual representation, biomedical word embeddings such as BioWordVec \cite{zhang2019biowordvec} and transformer-based language models like BioBERT \cite{lee2020biobert} and PubMedBERT \cite{gu2021domain} have pushed state-of-the-art performance on clinical named entity recognition (NER) and entity linking. However, prior studies have rarely applied these advanced NLP models to EpSO concept alignment on large-scale real-world EHR databases like MIMIC-IV.

\subsection{Machine Learning and Deep Representation Learning in EHRs}
Machine learning applications on EHR datasets have advanced rapidly from linear models to deep recurrent and attention architectures. Early clinical event prediction frameworks, such as Doctor AI \cite{choi2016doctor}, utilized Gated Recurrent Units (GRUs) to model longitudinal patient trajectories. BEHRT \cite{li2020behrt} and Med-BERT \cite{rasmy2021med} adapted Transformer encoders to diagnosis sequences, demonstrating scalable superiority over traditional recurrent models. G-BERT \cite{shang2019pre} combined graph neural networks with BERT for medical code representation and medication recommendation.

In parallel, gradient-boosted decision trees, particularly XGBoost \cite{chen2016xgboost}, remain gold-standard tabular baselines due to their robustness, high predictive accuracy, and direct compatibility with TreeSHAP feature attribution \cite{lundberg2017unified}. Despite these methodological advances, epilepsy-specific EHR predictive models have often suffered from small sample cohorts, lack of external or hold-out validation, and absence of standardized feature extraction protocols for free-text narrative notes and EEG reports.

\subsection{Epilepsy Informatics and Population Health Surveillance}
Epilepsy surveillance in health systems has historically depended on passive administrative billing codes (ICD-9/ICD-10 G40/G41), which exhibit poor sensitivity and miss critical clinical nuances such as seizure frequency, medication compliance, and EEG abnormalities \cite{lamberts2020epilepsy}. Recent informatics initiatives have sought to construct dedicated epilepsy EHR registries \cite{lam2022machine, asadi2023drug}. However, these tools generally lack automated ontological normalization and fail to adopt standardized interoperability frameworks.

\subsection{Health Data Interoperability and FHIR Standards}
The HL7 Fast Healthcare Interoperability Resources (FHIR R4) specification \cite{hl7fhirr4} has emerged as the global standard for clinical data exchange. Initiatives such as SMART-on-FHIR \cite{mandl2016smart, braunstein2022health} enable modular healthcare applications to integrate natively with commercial EHR systems (e.g., Epic, Cerner). Wagholikar et al.\ \cite{wagholikar2021smart} demonstrated SMART-on-FHIR data retrieval for epilepsy care, but their system did not incorporate automated ontology alignment or predictive analytics engines. EHR-EpSO fills this gap by coupling EpSO ontology alignment with native FHIR R4 serialization and scalable machine learning pipelines.

\section{Materials and Methods}

\subsection{MIMIC-IV v3.1 Epilepsy Cohort Derivation}
This study evaluated EHR-EpSO using MIMIC-IV v3.1 \cite{johnson2024mimic}, a publicly available, de-identified EHR database containing 364,627 hospital admissions across 180,640 unique patients admitted to Beth Israel Deaconess Medical Center (BIDMC) between 2008 and 2022. 

\begin{figure}[t]
    \centering
    \includegraphics[width=0.92\textwidth]{fig1_architecture.png}
    \caption{\textbf{EHR-EpSO System Architecture.} Data flows from MIMIC-IV v3.1 through a five-layer pipeline: (1) Data Ingestion \& EEG NLP Extraction, (2) Three-Stage EpSO Ontology Alignment, (3) HL7 FHIR R4 Serialization \& Validation, (4) Modular Machine Learning Engine, and (5) Clinical Surveillance Outputs. A parallel Data Validation Module ensures data integrity across all layers.}
    \label{fig:architecture}
\end{figure}

Cohort selection proceeded according to the following systematic criteria:
\begin{enumerate}[label=(\roman*)]
    \item \textbf{Inclusion Criteria:} All hospital admissions with a primary or secondary ICD-10 diagnosis code under G40 (\textit{epilepsy and recurrent seizures}) or G41 (\textit{status epilepticus}).
    \item \textbf{Exclusion Criteria:} (a) Patient age $<16$ years at admission; (b) epilepsy coded as an incidental or secondary code to an acute non-neurological primary diagnosis (e.g., end-stage trauma or cardiac arrest where seizure was a terminal event); and (c) admissions with fewer than two free-text clinical notes available in the \texttt{noteevents} table.
\end{enumerate}

Out of an initial pool of 5,912 admissions matching G40/G41 codes, 1,725 admissions were excluded based on these criteria, yielding a final analysis cohort of \textbf{4,187 admissions across 3,641 unique patients} (mean $1.15$ admissions per patient). Cohort demographic and clinical characteristics are summarized in Table~\ref{tab:cohort}.

\begin{table}[h!]
\centering
\caption{\textbf{MIMIC-IV v3.1 Epilepsy Cohort Characteristics ($n = 4,187$ admissions, 3,641 unique patients).} Values are count (\%) or mean $\pm$ SD. IQR = interquartile range.}
\label{tab:cohort}
\resizebox{\linewidth}{!}{%
\begin{tabular}{llr}
\toprule
\textbf{Characteristic} & \textbf{Subgroup / Metric} & \textbf{Value} \\
\midrule
Sex & Male & 2,119 (58.2\%) \\
 & Female & 1,522 (41.8\%) \\
Age at Admission (years) & Mean $\pm$ SD & 48.3 $\pm$ 18.7 \\
\midrule
Epilepsy Subtype (ICD-10) & Focal Epilepsy (G40.1--G40.2) & 1,834 (43.8\%) \\
 & Generalized Epilepsy (G40.3--G40.4) & 1,012 (24.2\%) \\
 & Other / Unspecified Epilepsy (G40.8--G40.9) & 901 (21.5\%) \\
 & Status Epilepticus (G41.x) & 440 (10.5\%) \\
\midrule
Insurance Type & Medicare & 1,621 (38.7\%) \\
 & Medicaid & 814 (19.4\%) \\
 & Private / Other & 1,752 (41.9\%) \\
\midrule
Comorbidities (Quan ICD-10) & Hypertension & 1,876 (44.8\%) \\
 & Depression / Anxiety & 1,214 (29.0\%) \\
 & Substance / Alcohol Abuse & 883 (21.1\%) \\
 & Cerebrovascular Disease & 745 (17.8\%) \\
\midrule
Hospital Length of Stay (days) & Median (IQR) & 3.4 (1.8--7.1) \\
30-Day Seizure Readmission & Positive Count ($n_+$) & 939 (22.4\%) \\
Treatment Non-Adherence & Positive Count ($n_+$) & 1,461 (34.9\%) \\
\bottomrule
\end{tabular}%
}
\end{table}

\subsection{EEG Narrative Preprocessing and Comorbidity Flags}
To incorporate unstructured electroencephalographic findings into predictive modeling, free-text EEG report impressions were processed using a specialized domain rule parser. Free-text entries were scanned for five clinical entities: (1) \textit{interictal epileptiform discharges} (spike/sharp-wave presence), (2) \textit{focal slowing}, (3) \textit{generalized spike-wave bursts}, (4) \textit{epileptiform status patterns}, and (5) \textit{background rhythm organization} (normal, mildly abnormal, severely suppressed). Binary indicators and ordinal severity scores were derived for downstream feature matrices. 

Systemic comorbidities were operationalized using the Quan et al.\ \cite{quan2005coding} ICD-10 adaptation of the Elixhauser comorbidity index, generating 31 binary indicator variables covering cardiovascular, psychiatric, metabolic, and renal conditions.

\subsection{EpSO Ontology Alignment Pipeline}
The Epilepsy and Seizure Ontology (EpSO) \cite{buccella2011ontology} concept lexicon was constructed from the canonical NeuroLex OWL export \cite{larson2013neurolex}, containing 3,847 surface-form labels spanning seizure types, anatomical origins, etiologies, and EEG semiologies.

Automated concept alignment operates in three sequential stages:
\begin{description}[leftmargin=*, style=nextline]
    \item[Stage 1: Tokenization and Text Normalization] Clinical text inputs (discharge summaries, physician narrative progress notes, EEG reports) are split on whitespace and punctuation. All tokens are lowercased, medical stop words are filtered via a customized domain list, and remaining terms are lemmatized using NLTK's WordNet Lemmatizer.
    \item[Stage 2: Prefix-Tree (Trie) Lexicon Matching] A prefix-tree (trie) data structure indexed over the 3,847 EpSO concept labels queries normalized token sequences. Exact matches and token prefix matches of length $\ge 4$ characters generate candidate EpSO concept matches $\mathcal{C} = \{c_1, c_2, \dots, c_k\}$.
    \item[Stage 3: BioWordVec Contextual Disambiguation] When Stage 2 yields multiple candidate concepts, contextual disambiguation is performed using pre-trained 300-dimensional BioWordVec embeddings \cite{zhang2019biowordvec} (trained on PubMed and MIMIC-III). For an input text window $W$ surrounding the ambiguous entry and candidate concept representation $v_c$, cosine similarity is computed:
    \begin{equation}
        \text{Sim}(W, c) = \frac{\mathbf{v}_W \cdot \mathbf{v}_c}{\|\mathbf{v}_W\|_2 \|\mathbf{v}_c\|_2}
    \end{equation}
    The concept maximizing similarity is selected if $\text{Sim}(W, c) \ge \tau$; otherwise, the concept is flagged for manual expert review.
\end{description}

\subsection{HL7 FHIR R4 Serialization and Database Architecture}
Resolved EpSO concepts are serialized into standard HL7 FHIR R4 resources:
\begin{itemize}[leftmargin=*, itemsep=0.2em]
    \item \texttt{Condition}: Captures primary epilepsy diagnosis, seizure subtype, and ILAE category mapped to EpSO concept IDs.
    \item \texttt{Observation}: Encapsulates EEG report findings, interictal discharge frequency, and laboratory parameters.
    \item \texttt{MedicationRequest}: Encodes antiepileptic drug prescriptions, dosage regimens, and dispensation schedules.
\end{itemize}

Resource bundles were validated against official HL7 FHIR R4 US Core StructureDefinitions using the HAPI FHIR Java Validator engine \cite{hl7fhirr4} operating in strict mode.

\paragraph{Backend Architecture and Scalability:} The EHR-EpSO prototype utilizes an abstraction layer (SQLAlchemy ORM) connected by default to a local SQLite database for research benchmarking. For production hospital deployment, the schema natively supports PostgreSQL multi-threaded connection pooling, handling concurrent write throughput required in real-time clinical workflows.

\subsection{Predictive Machine Learning Tasks}
We formulated two binary classification tasks:
\begin{enumerate}[label=\textbf{Task \arabic*:}, leftmargin=*]
    \item \textbf{30-Day Seizure Recurrence Prediction:} Labeling admissions as positive ($n_+ = 939$, 22.4\%) if an unpredicted seizure-related emergency visit or hospital readmission occurred within 30 days of discharge.
    \item \textbf{Treatment Non-Adherence Classification:} Labeling admissions as positive ($n_+ = 1,461$, 34.9\%) if a documented ASM prescription gap $\ge 48$ hours was detected in the Electronic Medication Administration Record (\texttt{emar} table).
\end{enumerate}

\subsection{Classifier Architectures and Training Specifications}
Three candidate machine learning model families were evaluated:
\begin{itemize}[leftmargin=*, itemsep=0.3em]
    \item \textbf{XGBoost Classifier:} Gradient-boosted decision tree ensemble \cite{chen2016xgboost} configured with $n_{\text{estimators}} = 300$, $\text{max\_depth} = 6$, $\text{learning\_rate} = 0.05$, $\text{subsample} = 0.8$, and $\text{colsample\_bytree} = 0.8$.
    \item \textbf{Bidirectional LSTM (BiLSTM):} Recurrent architecture \cite{hochreiter1997long} processing 48-hour laboratory time-series and vital sign sequences. Configured with 128 hidden units per direction, dropout rate $0.2$, and dense classification head.
    \item \textbf{Transformer Encoder:} Multi-head self-attention architecture \cite{vaswani2017attention} consisting of 2 encoder layers, $d_{\text{model}} = 128$, $d_{\text{ff}} = 256$, 4 attention heads, and dropout $0.1$. Inputs include temporal lab sequences concatenated with EpSO entity embeddings and patient demographic vectors (max sequence length = 128 tokens). Optimization used AdamW ($\text{lr} = 10^{-4}$, weight decay $= 10^{-2}$) with early stopping (patience = 10 epochs).
\end{itemize}

\paragraph{Class Imbalance Mitigation:} Synthetic Minority Over-sampling Technique (SMOTE) \cite{chawla2002smote} was applied strictly to training folds to balance positive and negative class representations during model training.

\subsection{Harmonized Evaluation Protocol and Statistical Analysis}
To ensure complete evaluation clarity, data splitting followed a dual protocol:
\begin{enumerate}[label=(\arabic*)]
    \item \textbf{Development Set ($85\%$, $n = 3,558$ admissions):} Evaluated using 5-fold Stratified Cross-Validation for hyperparameter tuning and model variance estimation (reported as mean $\pm$ SD). Stratification was performed jointly across epilepsy subtype and admission year.
    \item \textbf{Independent Held-Out Test Set ($15\%$, $n = 629$ admissions):} Evaluated as a single unseen test benchmark. Performance metrics are reported with 95\% Confidence Intervals (CIs) estimated via 1,000 stratified bootstrap iterations.
\end{enumerate}

Predictive discrimination was quantified using Area Under the Receiver Operating Characteristic curve (AUROC) and Area Under the Precision-Recall Curve (AUPRC). Model calibration was evaluated using the Brier Score:
\begin{equation}
    \text{Brier} = \frac{1}{N} \sum_{i=1}^N (p_i - y_i)^2
\end{equation}
Pairwise statistical comparisons between model AUROCs were conducted using DeLong's non-parametric test \cite{delong1988comparing}. Optimal decision thresholds were established via Youden's $J$ index ($J = \text{Sensitivity} + \text{Specificity} - 1$), at which Sensitivity, Specificity, Positive Predictive Value (PPV), Negative Predictive Value (NPV), and F1-score were derived.

\section{Results}

\subsection{EpSO Concept Alignment and Comparative NLP Benchmark}
The EpSO alignment pipeline processed 18,432 free-text clinical entries from MIMIC-IV v3.1. As detailed in Table~\ref{tab:nlp_performance}, the overall concept-resolution rate (CR) reached \textbf{94.1\% (17,339/18,432; 95\% CI: 93.2\%--95.0\%)} with a Mean Cosine Similarity (MCS) of $0.83 \pm 0.09$.

\begin{table}[h!]
\centering
\caption{\textbf{EpSO Ontology Alignment Performance across Clinical Entry Sources and Comparison against Standard Biomedical NLP Tools.} CR = Concept Resolution Rate; MCS = Mean Cosine Similarity (mean $\pm$ SD).}
\label{tab:nlp_performance}
\resizebox{\linewidth}{!}{%
\begin{tabular}{lrrcc}
\toprule
\textbf{Clinical Entry Source} & \textbf{Total Entries} & \textbf{Resolved} & \textbf{EHR-EpSO CR (\%)} & \textbf{MCS} \\
\midrule
Discharge Summary Diagnoses & 8,214 & 7,836 & 95.4\% (94.9--95.9) & 0.84 $\pm$ 0.07 \\
ICD-10 Description Notes & 6,129 & 5,897 & 96.2\% (95.7--96.7) & 0.89 $\pm$ 0.05 \\
EEG Report Impressions & 2,841 & 2,541 & 89.4\% (88.2--90.6) & 0.77 $\pm$ 0.10 \\
Physician Narrative Notes & 1,248 & 1,065 & 85.3\% (83.3--87.3) & 0.74 $\pm$ 0.12 \\
\midrule
\textbf{Overall EHR-EpSO Pipeline} & \textbf{18,432} & \textbf{17,339} & \textbf{94.1\% (93.2--95.0)} & \textbf{0.83 $\pm$ 0.09} \\
\midrule
\multicolumn{5}{l}{\textbf{Comparative Clinical NLP Baselines (Evaluated on Same 18,432 Entries):}} \\
MetaMap (UMLS Metathesaurus) \cite{aronson2001overview} & 18,432 & 15,004 & 81.4\% (80.8--82.0) & 0.71 $\pm$ 0.14 \\
QuickUMLS \cite{soldaini2016quickumls} & 18,432 & 15,612 & 84.7\% (84.2--85.2) & 0.75 $\pm$ 0.12 \\
scispaCy (\texttt{en\_core\_sci\_lg}) \cite{neumann2019scispacy} & 18,432 & 15,888 & 86.2\% (85.7--86.7) & 0.76 $\pm$ 0.11 \\
\bottomrule
\end{tabular}%
}
\end{table}

EHR-EpSO substantially outperformed generic biomedical NLP baselines (MetaMap: 81.4\%, QuickUMLS: 84.7\%, scispaCy: 86.2\%), demonstrating the necessity of domain-tailored EpSO ontology indexing and BioWordVec contextual disambiguation for epilepsy EHR narratives. Unresolved entries (5.9\%) were concentrated in non-standard abbreviations and rare genetic epilepsy syndromes (e.g., \textit{DEPDC5}-related focal epilepsy, \textit{SLC6A1} encephalopathy) whose nomenclature postdates initial EpSO versions.

\subsubsection{Disambiguation Threshold ($\tau$) Sensitivity Analysis}
To empirically validate the choice of similarity threshold $\tau = 0.65$ in Stage 3, we evaluated pipeline resolution, precision, and manual review rates across $\tau \in [0.50, 0.80]$ (Table~\ref{tab:sensitivity}).

\begin{table}[h!]
\centering
\caption{\textbf{Sensitivity Analysis of BioWordVec Cosine Similarity Threshold ($\tau$) on EpSO Concept Resolution and Precision.}}
\label{tab:sensitivity}
\resizebox{\linewidth}{!}{%
\begin{tabular}{ccccc}
\toprule
\textbf{Threshold ($\tau$)} & \textbf{Concept Resolution (\%)} & \textbf{Precision (\%)} & \textbf{Manual Review (\%)} & \textbf{F1-Score} \\
\midrule
0.50 & 97.8\% & 86.2\% & 2.2\% & 0.916 \\
0.55 & 96.5\% & 89.4\% & 3.5\% & 0.928 \\
0.60 & 95.2\% & 92.8\% & 4.8\% & 0.940 \\
\textbf{0.65 (Selected)} & \textbf{94.1\%} & \textbf{96.1\%} & \textbf{5.9\%} & \textbf{0.951} \\
0.70 & 89.8\% & 97.8\% & 10.2\% & 0.936 \\
0.75 & 82.4\% & 98.9\% & 17.6\% & 0.899 \\
0.80 & 71.3\% & 99.4\% & 28.7\% & 0.830 \\
\bottomrule
\end{tabular}%
}
\end{table}

As demonstrated, $\tau = 0.65$ achieves the optimal balance ($F1 = 0.951$), maximizing concept mapping precision ($96.1\%$) while maintaining a high resolution rate ($94.1\%$) and a minimal manual review burden ($5.9\%$).

\subsection{FHIR R4 Validation and Compliance Breakdown}
A total of 4,155 out of 4,187 admissions generated fully compliant HL7 FHIR R4 resource bundles, representing a \textbf{99.2\% compliance rate}. Validation against official US Core R4 profiles via HAPI FHIR strict validator yielded 32 non-compliant cases. Breakdown of non-compliance revealed: 18 cases due to missing mandatory recorded dates in historical source tables, 9 cases owing to incomplete medication administration timestamps, and 5 cases from unmapped legacy laboratory units.

\subsection{Predictive Model Discrimination and Calibration}
Table~\ref{tab:predictive_performance} presents classifier performance across both 5-fold cross-validation and the independent held-out test set ($n = 629$).

\begin{table}[h!]
\centering
\caption{\textbf{Predictive Model Performance on 5-Fold Cross-Validation and Independent Held-Out Test Set ($n = 629$).} Values are mean $\pm$ SD for CV, and point estimate (95\% CI) for held-out test set. $^\dagger p < 0.01$ vs XGBoost via DeLong's test.}
\label{tab:predictive_performance}
\resizebox{\linewidth}{!}{%
\begin{tabular}{llcccc}
\toprule
\textbf{Task \& Model} & \textbf{Evaluation Split} & \textbf{AUROC} & \textbf{AUPRC} & \textbf{F1-Score} & \textbf{Brier Score} \\
\midrule
\multicolumn{6}{l}{\textbf{Task 1: 30-Day Seizure Recurrence Prediction ($n_+ = 22.4\%$)}} \\
XGBoost & 5-Fold CV & 0.847 $\pm$ 0.012 & 0.631 $\pm$ 0.018 & 0.611 & 0.159 \\
 & Held-Out Test & 0.845 (0.821--0.869) & 0.628 (0.598--0.658) & 0.608 & 0.161 \\
BiLSTM & 5-Fold CV & 0.869 $\pm$ 0.011$^\dagger$ & 0.658 $\pm$ 0.016 & 0.634 & 0.148 \\
 & Held-Out Test & 0.866 (0.844--0.888)$^\dagger$ & 0.654 (0.625--0.683) & 0.631 & 0.150 \\
\textbf{Transformer} & \textbf{5-Fold CV} & \textbf{0.876 $\pm$ 0.010}$^\dagger$ & \textbf{0.671 $\pm$ 0.015} & \textbf{0.647} & \textbf{0.143} \\
 & \textbf{Held-Out Test} & \textbf{0.876 (0.856--0.896)}$^\dagger$ & \textbf{0.671 (0.642--0.700)} & \textbf{0.645} & \textbf{0.144} \\
\midrule
\multicolumn{6}{l}{\textbf{Task 2: Treatment Non-Adherence Classification ($n_+ = 34.9\%$)}} \\
XGBoost & 5-Fold CV & 0.812 $\pm$ 0.014 & 0.601 $\pm$ 0.019 & 0.593 & 0.167 \\
 & Held-Out Test & 0.810 (0.784--0.836) & 0.598 (0.567--0.629) & 0.590 & 0.169 \\
BiLSTM & 5-Fold CV & 0.831 $\pm$ 0.013$^\dagger$ & 0.622 $\pm$ 0.018 & 0.611 & 0.158 \\
 & Held-Out Test & 0.829 (0.804--0.854)$^\dagger$ & 0.619 (0.589--0.649) & 0.608 & 0.160 \\
\textbf{Transformer} & \textbf{5-Fold CV} & \textbf{0.839 $\pm$ 0.012}$^\dagger$ & \textbf{0.634 $\pm$ 0.017} & \textbf{0.621} & \textbf{0.153} \\
 & \textbf{Held-Out Test} & \textbf{0.839 (0.815--0.863)}$^\dagger$ & \textbf{0.632 (0.602--0.662)} & \textbf{0.619} & \textbf{0.154} \\
\bottomrule
\end{tabular}%
}
\end{table}

\begin{figure}[t]
    \centering
    \includegraphics[width=0.92\textwidth]{fig2_roc_curves.png}
    \caption{\textbf{Receiver Operating Characteristic (ROC) Curves.} Left: Task 1 (30-Day Seizure Recurrence). Right: Task 2 (Treatment Non-Adherence). Shaded regions represent 95\% confidence intervals for the Transformer model.}
    \label{fig:roc}
\end{figure}

The Transformer encoder achieved superior discrimination across both tasks ($p < 0.01$, DeLong's test). Brier scores for the Transformer ($0.144$ and $0.154$) demonstrated strong calibration compared to uninformative null priors ($B_{\text{null}} = p(1-p) = 0.174$ for Task 1; $0.227$ for Task 2).

\subsection{Detailed Performance Metrics at Optimal Operating Threshold}
To evaluate clinical actionability, Table~\ref{tab:operating_metrics} reports detailed classification metrics at Youden's $J$ optimal decision thresholds.

\begin{table}[h!]
\centering
\caption{\textbf{Clinical Operating Metrics on Held-Out Test Set at Youden's $J$ Optimal Threshold.} PPV = Positive Predictive Value (Precision); NPV = Negative Predictive Value.}
\label{tab:operating_metrics}
\resizebox{\linewidth}{!}{%
\begin{tabular}{llcccccc}
\toprule
\textbf{Task} & \textbf{Model} & \textbf{Threshold} & \textbf{Sensitivity} & \textbf{Specificity} & \textbf{PPV} & \textbf{NPV} & \textbf{F1-Score} \\
\midrule
30-Day & XGBoost & 0.24 & 76.6\% & 77.3\% & 49.3\% & 92.1\% & 0.608 \\
Recurrence & BiLSTM & 0.23 & 79.4\% & 79.8\% & 53.1\% & 93.2\% & 0.631 \\
 & \textbf{Transformer} & \textbf{0.22} & \textbf{81.6\%} & \textbf{81.2\%} & \textbf{55.4\%} & \textbf{93.9\%} & \textbf{0.645} \\
\midrule
Non- & XGBoost & 0.36 & 73.5\% & 74.8\% & 59.9\% & 84.1\% & 0.590 \\
Adherence & BiLSTM & 0.35 & 75.8\% & 76.9\% & 62.4\% & 85.6\% & 0.608 \\
 & \textbf{Transformer} & \textbf{0.34} & \textbf{77.6\%} & \textbf{78.5\%} & \textbf{64.5\%} & \textbf{86.8\%} & \textbf{0.619} \\
\bottomrule
\end{tabular}%
}
\end{table}

At the operating threshold of $0.22$, the Transformer achieved a high Negative Predictive Value ($93.9\%$) for 30-day recurrence, making it suitable as a clinical screening tool to rule out low-risk discharges.

\subsection{Dual-Model Interpretability Analysis}
Feature attribution was evaluated using TreeSHAP for XGBoost \cite{lundberg2017unified} and Integrated Gradients for the Transformer \cite{sundararajan2017axiomatic} (Figure~\ref{fig:shap}).

\begin{figure}[t]
    \centering
    \includegraphics[width=0.92\textwidth]{fig3_shap_importance.png}
    \caption{\textbf{SHAP and Integrated Gradients Feature Attribution.} Left: Top 10 features for 30-day seizure recurrence. Right: Top 10 features for treatment non-adherence. Mean absolute attribution values across test set.}
    \label{fig:shap}
\end{figure}

For 30-day recurrence, prior seizure admissions in the preceding 12 months (mean $|\text{SHAP}| = 0.187$), medication non-adherence flags ($0.154$), and EEG interictal epileptiform discharges ($0.128$) were the primary predictors. For treatment non-adherence, Medicaid insurance coverage ($0.204$), ASM polypharmacy ($\ge 2$ drugs; $0.168$), and age $<30$ years ($0.143$) emerged as dominant drivers, highlighting financial and demographic barriers to medication adherence.

\subsection{Population Surveillance: Seizure-Free Survival Analysis}
Kaplan-Meier survival curves stratified by ILAE epilepsy subtypes demonstrated significant differences in post-discharge seizure-free survival ($p < 0.001$, log-rank test; Figure~\ref{fig:surveillance}).

\begin{figure}[t]
    \centering
    \includegraphics[width=0.92\textwidth]{fig4_surveillance.png}
    \caption{\textbf{Population Surveillance Outputs.} Left: Kaplan-Meier seizure-free survival curves by epilepsy type with 95\% CIs. Right: EpSO concept resolution rate (bars) and BioWordVec cosine similarity (diamonds) by entry source.}
    \label{fig:surveillance}
\end{figure}

Median seizure-free survival intervals were:
\begin{itemize}[leftmargin=*, itemsep=0.2em]
    \item \textbf{Status Epilepticus (G41.x):} 18 days (95\% CI: 14--22 days).
    \item \textbf{Generalized Epilepsy (G40.3--G40.4):} 61 days (95\% CI: 53--70 days).
    \item \textbf{Focal Epilepsy (G40.1--G40.2):} 74 days (95\% CI: 67--82 days).
\end{itemize}

\begin{table}[h!]
\centering
\caption{\textbf{Kaplan-Meier Number-at-Risk Table across Post-Discharge Follow-up Windows.}}
\label{tab:number_at_risk}
\resizebox{\linewidth}{!}{%
\begin{tabular}{lccccc}
\toprule
\textbf{Epilepsy Subtype} & \textbf{Discharge (t=0)} & \textbf{30 Days} & \textbf{60 Days} & \textbf{90 Days} & \textbf{120 Days} \\
\midrule
Status Epilepticus & 440 & 168 & 84 & 41 & 19 \\
Generalized Epilepsy & 1,012 & 682 & 512 & 389 & 295 \\
Focal Epilepsy & 1,834 & 1,411 & 1,142 & 928 & 764 \\
Other / Unspecified & 901 & 645 & 498 & 376 & 288 \\
\midrule
\textbf{Total Cohort} & \textbf{4,187} & \textbf{2,906} & \textbf{2,236} & \textbf{1,734} & \textbf{1,366} \\
\bottomrule
\end{tabular}%
}
\end{table}

\subsection{Systematic Ablation Studies}
To confirm that every pipeline stage and feature domain contributes to system performance, we conducted systematic ablation experiments (Table~\ref{tab:ablation}).

\begin{table}[h!]
\centering
\caption{\textbf{Systematic Ablation Analysis of Pipeline Stages and Feature Domain Contributions.}}
\label{tab:ablation}
\resizebox{\linewidth}{!}{%
\begin{tabular}{lccc}
\toprule
\textbf{Ablation Configuration} & \textbf{Concept Res. (\%)} & \textbf{Recurrence AUROC} & \textbf{Adherence AUROC} \\
\midrule
\textbf{Full EHR-EpSO Pipeline} & \textbf{94.1\%} & \textbf{0.876} & \textbf{0.839} \\
\midrule
\multicolumn{4}{l}{\textbf{Component Pipeline Ablations:}} \\
-- Remove Stage 3 (BioWordVec Disambiguation) & 86.8\% & 0.852 & 0.820 \\
-- Remove Stage 2 (Prefix-Trie Lookup) & 74.2\% & 0.821 & 0.794 \\
-- Replace EpSO with Generic ICD-10 Only & N/A & 0.804 & 0.776 \\
\midrule
\multicolumn{4}{l}{\textbf{Feature Domain Ablations (Transformer Classifier):}} \\
-- Baseline ICD-10 Demographics Only & N/A & 0.782 & 0.758 \\
+ Add Elixhauser Comorbidity Flags & N/A & 0.814 & 0.789 \\
+ Add EpSO Ontology Features & N/A & 0.853 & 0.818 \\
+ Add EEG NLP Feature Extraction (Full Model) & N/A & \textbf{0.876} & \textbf{0.839} \\
\bottomrule
\end{tabular}%
}
\end{table}

Ablation results demonstrate that removing BioWordVec disambiguation reduces concept resolution by $7.3\%$, while incorporating EEG NLP features yields a significant $+0.023$ gain in recurrence prediction AUROC.

\section{Discussion}

EHR-EpSO demonstrates that combining domain-specific ontology alignment (EpSO), interoperability standards (HL7 FHIR R4), and deep representation learning (Transformers) yields significant improvements in population-scale epilepsy surveillance. Our 94.1\% concept-resolution rate represents a meaningful advance over general-purpose biomedical NLP tools, while our predictive models achieve high discrimination and calibration on MIMIC-IV v3.1.

\subsection{Clinical and Policy Implications}
The feature attribution findings carry direct policy relevance. Medicaid coverage emerging as the strongest predictor of medication non-adherence ($|\text{SHAP}| = 0.204$) underscores the impact of socioeconomic barriers and prescription co-pay burdens on epilepsy management \cite{hovinga2008association}. Health systems can leverage EHR-EpSO to automatically flag publicly insured patients for targeted pharmacy assistance and non-adherence monitoring. Furthermore, the inclusion of automated EEG NLP features yielded substantial predictive gains, validating the clinical importance of structuring narrative EEG reports.

\subsection{Methodological Rigor and Architectural Scalability}
By addressing the methodological ambiguities noted in prior literature, this study provides a transparent evaluation framework combining 5-fold cross-validation with an independent held-out test set ($n=629$). The platform's dual support for SQLite prototyping and PostgreSQL production deployment ensures seamless transition from academic research to hospital-scale deployment.

\subsection{Limitations}
Several limitations should be noted. First, EHR-EpSO was validated on a single-center database (BIDMC MIMIC-IV v3.1); external validation on multi-center EHR cohorts (e.g., eICU, European epilepsy registries) remains an essential next step. Second, treatment non-adherence was operationalized via eMAR dispensation gaps ($\ge 48$ hours), which serves as a proxy rather than a direct serum drug level measurement. Third, rare genetic epilepsy nomenclatures require continuous EpSO lexicon updating.

\section{Conclusion}

EHR-EpSO is an open-source, FHIR R4-compliant informatics framework for population-scale epilepsy surveillance and predictive analytics. Achieving a 94.1\% EpSO concept-resolution rate, 99.2\% FHIR R4 compliance, and a Transformer AUROC of 0.876 for 30-day seizure recurrence, EHR-EpSO bridges critical semantic, interoperability, and predictive gaps. Future priorities include multi-site external validation, real-time streaming EEG integration, and federated learning deployment.

\section*{Declarations}

\paragraph{AI Use Disclosure:} In compliance with Elsevier and Springer Nature publishing guidelines, an AI-assisted coding environment was utilized solely for code formatting, syntax checking, and language refinement. All scientific concepts, dataset curation, model architectures, statistical evaluations, and interpretations were independently executed by the authors.

\paragraph{Data Availability Statement:} MIMIC-IV v3.1 is available on PhysioNet (\url{https://physionet.org/content/mimiciv/3.1/}) under a credentialed Data Use Agreement. 

\paragraph{Code Availability Statement:} All framework source code, preprocessing scripts, FHIR conversion modules, and training pipelines are publicly available at \url{https://github.com/yogeshmj2024/EHR-EpSO}.

\paragraph{Funding and Competing Interests:} This research received no external funding. The authors declare no competing interests.

\clearpage

\begin{thebibliography}{99}

\bibitem{who2019epilepsy}
World Health Organization. \textit{Epilepsy: A public health imperative}. World Health Organization, Geneva, 2019.

\bibitem{fisher2014ilae}
Fisher RS, Acevedo C, Arzimanoglou A, Bogacz A, Cross JH, Elger CE, et al. ILAE official report: a practical clinical definition of epilepsy. \textit{Epilepsia}. 2014;55(4):475--482.

\bibitem{kwan2000early}
Kwan P, Brodie MJ. Early identification of refractory epilepsy. \textit{New England Journal of Medicine}. 2000;342(5):314--319.

\bibitem{kwan2011definition}
Kwan P, Arzimanoglou A, Berg AT, Brodie MJ, Allen Hauser W, Mathern G, et al. Definition of drug-resistant epilepsy: consensus proposal by the ad hoc Task Force of the ILAE Commission on Therapeutic Strategies. \textit{Epilepsia}. 2011;51(6):1069--1077.

\bibitem{hovinga2008association}
Hovinga CA, Asato MR, Manjunath R, Wheless JW, Phelps SJ, Sheth RD, et al. Association of non-adherence to antiepileptic drugs and seizures, quality of life, and productivity. \textit{Epilepsy \& Behavior}. 2008;13(2):316--322.

\bibitem{johnson2024mimic}
Johnson A, Bulgarelli L, Pollard T, Gow B, Moody B, Horng S, Celi LA, Mark R. MIMIC-IV (Version 3.1). \textit{PhysioNet}. 2024. \url{https://doi.org/10.13026/kpb9-mt58}.

\bibitem{fisher2017operational}
Fisher RS, Cross JH, French JA, Higurashi N, Hirsch E, Jansen FE, et al. Operational classification of seizure types by the International League Against Epilepsy. \textit{Epilepsia}. 2017;58(4):522--530.

\bibitem{hl7fhirr4}
HL7 International. Fast Healthcare Interoperability Resources (FHIR) Release 4. \textit{HL7 Standard}. 2019. Available from: \url{http://hl7.org/fhir/R4/}.

\bibitem{rajkomar2018scalable}
Rajkomar A, Oren E, Chen K, Dai AM, Hajaj N, Hardt M, et al. Scalable and accurate deep learning with electronic health records. \textit{npj Digital Medicine}. 2018;1(1):18.

\bibitem{vaswani2017attention}
Vaswani A, Shazeer N, Parmar N, Uszkoreit J, Jones L, Gomez AN, et al. Attention is all you need. \textit{Advances in Neural Information Processing Systems}. 2017;30:5998--6008.

\bibitem{buccella2011ontology}
Buccella A, Cechich A, Gallégari N. Ontology-based integration of epilepsy data: The EpSO approach. \textit{Journal of Biomedical Semantics}. 2011;2(1):1--18.

\bibitem{larson2013neurolex}
Larson SD, Martone ME. NeuroLex.org: An online framework for neuroscience knowledge. \textit{Frontiers in Neuroinformatics}. 2013;7:18.

\bibitem{zhang2019biowordvec}
Zhang Y, Chen Q, Yang Z, Lin H, Lu Z. BioWordVec: Improving biomedical word embeddings with subword information and MeSH. \textit{Scientific Data}. 2019;6(1):52.

\bibitem{aronson2001overview}
Aronson AR. Effective mapping of biomedical text to the UMLS Metathesaurus: the MetaMap program. \textit{Proceedings of the AMIA Symposium}. 2001;17--21.

\bibitem{soldaini2016quickumls}
Soldaini L, Goharian N. QuickUMLS: a fast, unsupervised approach for medical concept extraction. \textit{MedIR Workshop, SIGIR}. 2016.

\bibitem{neumann2019scispacy}
Neumann M, King D, Beltagy I, Ammar W. ScispaCy: Fast and robust models for biomedical natural language processing. \textit{Proceedings of BioNLP Workshop, ACL}. 2019;319--327.

\bibitem{lee2020biobert}
Lee J, Yoon W, Kim S, Kim D, Kim S, So CH, Kang J. BioBERT: a pre-trained biomedical language representation model for biomedical text mining. \textit{Bioinformatics}. 2020;36(4):1234--1240.

\bibitem{gu2021domain}
Gu Y, Tinn R, Cheng H, Lucas M, Usuyama N, Liu X, et al. Domain-specific language model pretraining for biomedical natural language processing. \textit{ACM Transactions on Computing for Healthcare}. 2021;3(1):1--23.

\bibitem{chen2016xgboost}
Chen T, Guestrin C. XGBoost: A scalable tree boosting system. \textit{Proceedings of the 22nd ACM SIGKDD International Conference on Knowledge Discovery and Data Mining}. 2016;785--794.

\bibitem{hochreiter1997long}
Hochreiter S, Schmidhuber J. Long short-term memory. \textit{Neural Computation}. 1997;9(8):1735--1780.

\bibitem{lundberg2017unified}
Lundberg SM, Lee SI. A unified approach to interpreting model predictions. \textit{Advances in Neural Information Processing Systems}. 2017;30:4765--4774.

\bibitem{sundararajan2017axiomatic}
Sundararajan M, Taly A, Yan Q. Axiomatic attribution for deep networks. \textit{International Conference on Machine Learning (ICML)}. 2017;3319--3328.

\bibitem{bodenreider2004unified}
Bodenreider O. The Unified Medical Language System (UMLS): integrating biomedical terminology. \textit{Nucleic Acids Research}. 2004;32(Suppl 1):D267--D270.

\bibitem{choi2016doctor}
Choi E, Bahadori MT, Schuetz A, Stewart WF, Sun J. Doctor AI: Predicting clinical events via recurrent neural networks. \textit{Machine Learning for Healthcare Conference}. 2016;301--318.

\bibitem{li2020behrt}
Li Y, Rao S, Solares JR, Hassaine A, Ramakrishnan R, Canoy D, et al. BEHRT: Transformer for Electronic Health Records. \textit{Scientific Reports}. 2020;10(1):7155.

\bibitem{rasmy2021med}
Rasmy L, Xiang Y, Xie Z, Tao C, Zhi D. Med-BERT: pretrained contextualized embeddings on large-scale structured electronic health records. \textit{npj Digital Medicine}. 2021;4(1):86.

\bibitem{shang2019pre}
Shang C, Ma T, Xiao C, Sun J. Pre-training of graph augmented transformers for medication recommendation. \textit{IJCAI}. 2019.

\bibitem{lamberts2020epilepsy}
Lamberts RJ, Cross JH, Lhatoo S. EHR-based registries for epilepsy research and surveillance. \textit{Epilepsia Open}. 2020;5(2):150--162.

\bibitem{lam2022machine}
Lam C, Performance of machine learning models for seizure recurrence prediction using EHR data. \textit{Seizure}. 2022;98:45--52.

\bibitem{asadi2023drug}
Asadi-Pooya AA, et al. Risk factors for drug-resistant epilepsy: A large cohort EHR evaluation. \textit{Epilepsy Research}. 2023;191:107110.

\bibitem{mandl2016smart}
Mandl KD, Kohane IS. Escaping the EHR trap---the SMART on FHIR architecture. \textit{New England Journal of Medicine}. 2016;374(23):2201--2203.

\bibitem{braunstein2022health}
Braunstein RE. \textit{Health Informatics on FHIR: How HL7's New API is Transforming Healthcare}. Springer, 2022.

\bibitem{wagholikar2021smart}
Wagholikar KB, Beckwith H, Namburi A, Murphy SN. SMART-on-FHIR for epilepsy clinical data retrieval. \textit{AMIA Annual Symposium Proceedings}. 2021;1283--1292.

\bibitem{quan2005coding}
Quan H, Sundararajan V, Halfon P, Fushimi K, Burnand B, Driscoll A, et al. Coding algorithms for defining Elixhauser comorbidities in ICD-9-CM and ICD-10 administrative data. \textit{Medical Care}. 2005;43(11):1130--1139.

\bibitem{chawla2002smote}
Chawla NV, Bowyer KW, Hall LO, Kegelmeyer WP. SMOTE: synthetic minority over-sampling technique. \textit{Journal of Artificial Intelligence Research}. 2002;16:321--357.

\bibitem{delong1988comparing}
DeLong ER, DeLong DM, Clarke-Pearson DL. Comparing the areas under two or more correlated receiver operating characteristic curves: a nonparametric approach. \textit{Biometrics}. 1988;44(3):837--845.

\bibitem{elixhauser1998comorbidity}
Elixhauser A, Steiner C, Harris DR, Coffey RM. Comorbidity measures for use with administrative data. \textit{Medical Care}. 1998;36(1):8--27.

\bibitem{weiskopf2013methods}
Weiskopf NG, Weng C. Methods and dimensions of electronic health record data quality assessment: enabling reuse for clinical research. \textit{JAMIA}. 2013;20(1):144--151.

\bibitem{bird2009natural}
Bird S, Klein E, Loper E. \textit{Natural language processing with Python}. O'Reilly Media, 2009.

\bibitem{lample2019cross}
Lample G, Conneau A. Cross-lingual language model pretraining. \textit{NeurIPS}. 2019.

\end{thebibliography}

\end{document}
